\section{Completion of the induction}
\label{sec:completion}

We now combine the two band quotients with the dominance case. The
combinatorial point is that a partition outside the dominance range either
has a long enough first row for a width-one cut, or has enough boxes in its
first four rows for a width-two cut.

\begin{samepage}
\begin{lemma}[Horizontal-strip reduction]
\label{comb:strip-reduction}
Let $m\ge2$ and $\lambda\vdash N_m$. If $\lambda$ is not dominated by
$\rho_m$, then at least one of the following holds.
\begin{enumerate}[label=\textup{(\roman*)}]
\item There is a partition $\nu\vdash N_{m-1}$ such that
      $\lambda/\nu$ is a horizontal strip of size $m$.
\item One has $m\ge5$, and there are partitions
      \[
      \nu=\lambda^{(0)}\subset\lambda^{(1)}\subset\cdots
      \subset\lambda^{(q)}=\lambda,
      \qquad q\le4,
      \]
      with $\nu\vdash N_{m-2}$ and every
      $\lambda^{(i)}/\lambda^{(i-1)}$ a nonempty horizontal strip.
      Their total size is $2m-1$.
\end{enumerate}
\end{lemma}
\end{samepage}

\begin{proof}
Unlike the deletion in \Cref{lem:young-rule}, this reduction need not
preserve a prescribed dominance inequality; a rightmost suffix suffices
to preserve the partition.
If $\lambda_1\ge m$, remove the bottom box in each of the rightmost $m$
nonempty columns. Subtracting one from this suffix of the column-height
sequence preserves weak decrease and nonnegativity. The remaining boxes
therefore form a partition $\nu$, and the deleted boxes form a horizontal
strip. Since $N_m-m=N_{m-1}$, this proves \textup{(i)}.

Suppose instead that $\lambda_1\le m-1$. Failure of domination gives an
index $k<m$ with
\begin{equation}
\label{comb:dominance-failure}
\sum_{i=1}^k\lambda_i>km-\frac{k(k-1)}2.
\end{equation}
The index is less than $m$ because the first $m$ parts of the staircase
already have sum $N_m$. For $k\le3$, however,
\[
\sum_{i=1}^k\lambda_i\le k(m-1)
\le km-\frac{k(k-1)}2.
\]
Thus $4\le k<m$, and in particular $m\ge5$. Because the parts decrease,
the mean of the first four is at least the mean of the first $k$. Hence
\begin{equation}
\label{comb:four-row-capacity}
\sum_{i=1}^4\lambda_i
\ge\frac4k\sum_{i=1}^k\lambda_i
>4m-2(k-1)\ge2m+4>2m-1.
\end{equation}

A full sweep removes the bottom box of every nonempty column. If the
current row lengths are $(\alpha_1,\alpha_2,\ldots)$, subtracting one from
all its positive column heights gives row lengths
$(\alpha_2,\alpha_3,\ldots)$. The remaining diagram is nested in the
original one, and the removed cells form a horizontal strip of size
$\alpha_1$. Four full sweeps starting from $\lambda$ would therefore
remove $\lambda_1+\cdots+\lambda_4$ boxes.

Perform full sweeps until the remaining number of boxes to be removed is
at most the number of nonempty columns, then remove that number of bottom
boxes from the rightmost columns. The suffix argument in the first
paragraph shows that this last partial sweep also leaves a partition.
By \Cref{comb:four-row-capacity}, at most four nonempty sweeps remove
exactly $2m-1$ boxes. Their remainder has size
$N_m-(2m-1)=N_{m-2}$. Reversing the sequence gives \textup{(ii)}.
\end{proof}

Four strips can be necessary. Take $m=13$ and $\lambda=(8^{11},3)$,
where the exponent denotes eleven repeated parts equal to $8$.
This partition has size $91=N_{13}$ and fails domination at $k=12$,
since its first twelve parts sum to $91>90$.
Every horizontal strip in its eight-column diagram has at most eight
boxes, so three strips cannot remove the required $25=2m-1$ boxes.
The sweeps remove $8+8+8+1=25$ boxes and leave
$\nu=(8^7,7,3)$ of size $66=N_{11}$.

\begin{proof}[Proof of \Cref{thm:cyclic-saxl}]
We induct on $m$. The module $W_1$ is the one-dimensional representation
of the one-element group, so the assertion holds at $m=1$.

Let $m\ge2$, suppose the assertion holds for all smaller positive indices,
and fix $\lambda\vdash N_m$. If $\rho_m\unrhd\lambda$, then
\Cref{prop:dominated} supplies $S^\lambda$ in $W_m$. Otherwise apply
\Cref{comb:strip-reduction}.

In case \textup{(i)}, use the band cut with $s=1$, so
$|D|=N_{m-1}$ and $|F|=m$. The width-one band module $U_{m,1}$ is
trivial. By the induction hypothesis, $S_D^\nu$ is a constituent of
$W_{m-1}$. By \Cref{lem:pieri}, $S_{B_m}^\lambda$ occurs in
\[
\Ind_{S_D\times S_F}^{S_{B_m}}
\bigl(S_D^\nu\boxtimes\mathbf1\bigr).
\]
Consequently it occurs in the quotient from \Cref{prop:band-quotient}
and hence in $W_m$.

In case \textup{(ii)}, put $K=2m-1$ and use the cut with $s=2$.
Write $b_i=|\lambda^{(i)}|-|\lambda^{(i-1)}|$, so that
$b_1+\cdots+b_q=K$ and $q\le4$. Choose disjoint subsets of $F$ of
these sizes and form the permutation module
\[
A=\Ind_{S_{b_1}\times\cdots\times S_{b_q}}^{S_F}\mathbf1
  \cong\bigoplus_{\eta\vdash K}(S_F^\eta)^{\oplus a_\eta}.
\]
Repeated application of \Cref{lem:pieri}, followed by transitivity of
induction, shows that $S_{B_m}^\lambda$ is a constituent of
$\Ind_{S_D\times S_F}^{S_{B_m}}(S_D^\nu\boxtimes A)$. Equivalently,
\begin{equation}
\label{comb:positive-multiplicity}
\begin{gathered}
\sum_{\eta\vdash K}a_\eta c_{\nu,\eta}^{\lambda}>0,\\[3pt]
c_{\nu,\eta}^{\lambda}
=\dim\Hom_{S_{B_m}}\!\left(
 S_{B_m}^\lambda,
 \Ind_{S_D\times S_F}^{S_{B_m}}
 (S_D^\nu\boxtimes S_F^\eta)\right).
\end{gathered}
\end{equation}
Here the coefficients are nonnegative integers. Moreover, every
constituent of $A$ has at most $q$ rows: starting from the empty diagram,
each of the $q$ horizontal-strip additions can introduce at most one new
row, because two new rows would both require a new box in the first
column. Thus \Cref{comb:positive-multiplicity} provides some $\eta$ with
at most four rows, $a_\eta>0$, and $c_{\nu,\eta}^{\lambda}>0$.

\begin{samepage}
By \Cref{prop:two-band}, this $S_F^\eta$ occurs in $U_{m,2}$.
By induction, $S_D^\nu$ occurs in the specified cyclic module
$W_{m-2}$. One copy of each suffices to give $S_{B_m}^\lambda$ in
\[
\Ind_{S_D\times S_F}^{S_{B_m}}(W_{m-2}\boxtimes U_{m,2}),
\]
which is a quotient of $W_m$ by \Cref{prop:band-quotient}.
Complete reducibility again gives the constituent in $W_m$ itself.
This completes the induction. Notice that case \textup{(ii)} has
$m\ge5$, so every smaller module invoked has a positive index.
\end{samepage}
\end{proof}

\begin{proof}[Proof of \Cref{thm:saxl}]
The two factors generating $W_m$ have Specht type $\rho_m$.
Thus $W_m$ is a submodule of
$S^{\rho_m}\otimes_{\C}S^{\rho_m}$ under the diagonal action.
\Cref{thm:cyclic-saxl} supplies every $S^\lambda$ in that submodule,
so $g(\rho_m,\rho_m,\lambda)>0$ for every permitted $m$ and $\lambda$.
\end{proof}
